\documentclass[11pt,a4paper]{article}
\usepackage[margin=1in]{geometry}
\usepackage[T1]{fontenc}
\usepackage{lmodern}
\usepackage{microtype}
\usepackage{hyperref}
\usepackage{array}
\usepackage{longtable}
\hypersetup{
  colorlinks=true,
  linkcolor=black,
  urlcolor=black,
  pdftitle={OC Core 1.3.2 Readable Overview},
  pdfauthor={Alexander Yashin},
  pdfsubject={Ontology of Continua Core 1.3.2 readable overview},
  pdfkeywords={Ontology of Continua, OC Core 1.3.2, overview}
}
\title{\textbf{OC Core 1.3.2 Readable Overview}\\
\large What the Release Contains, What It Supports, and What It Deliberately Leaves Outside}
\author{Alexander Yashin\\Independent Researcher, Leipzig/Halle, Germany\\ORCID 0009-0008-6166-0914}
\date{Version v1.3.2; release date: 28.04.2026}
\begin{document}
\maketitle

\section*{Purpose}
OC Core 1.3.2 is a full no-send release-candidate package for owner review. Its role is to publish a stable, evidence-scoped baseline of the Core 1.3 line once owner and publication locks are lawfully released. It imports the public-safe 1.4 formal and benchmark elements that already have proof, replay, or gate support, while keeping private drafts and IP-sensitive material out of the public package.

\section*{Plain-Language Summary}
The Ontology of Continua studies objects that persist only while their admissible realizations remain live. When all admissible realizations are gone, the original continuum is dead. Something may remain, but that remainder is residue. If a later live continuum grows from that residue, OC Core classifies it as rebirth, not as persistence of the original object.

\section*{What Is New in 1.3.2}
This release integrates the current validated science into a coherent package:
\begin{itemize}
  \item a rebuilt master monograph with current 1.3.2 identity and Maria dedication;
  \item a release terminality ledger for the eighty-two tracked scientific release blockers;
  \item K0--K12 projection/replay state and domain packets;
  \item DRT strict-salvage state;
  \item Tsukrov evidence-first repair and reviewer navigation;
  \item Parfitian Cerberus pass for the current no-send candidate;
  \item LRGEF package governance, provenance, checksums, and no-send lock;
  \item night-delta integration represented as bounded patch-plan groups rather than private transcripts.
\end{itemize}

\section*{What Counts as Ready Science Here}
Ready science means that the claim is supported by proof/evidence, reproducible replay, source audit, or a release-visible terminal owner gate. A task note by itself is not science. A generic summary is not science. A strong phrase such as ``prediction'' or ``theorem'' is allowed only when the required support route exists in the release package.

\section*{What Is Outside the Release}
OC Core 1.4 race-hardening, private IP triage, private reviewer-source material, internal model transcripts, account or billing material, and unfinished research programs remain outside this public-safe package. They can be referenced only as future-work or owner-review context when the reference is already public-safe.

\section*{Publication Posture}
The package is prepared as \textbf{RELEASE\_READY\_NO\_SEND}. That means the release machine can build and evaluate it locally, but public push, publication, DOI minting, Zenodo upload, GitHub release creation, and outbound communication remain false until owner approval and final publication checks.

\section*{How to Read the Package}
\begin{longtable}{>{\raggedright\arraybackslash}p{0.30\linewidth}>{\raggedright\arraybackslash}p{0.62\linewidth}}
\textbf{Question} & \textbf{Release surface}\\
\hline
What is the theory? & Master monograph and journal core.\\
What exactly is supported? & Claim ledger, theorem roadmap, proof machinery appendix, and evidence matrix.\\
Were the blockers handled? & Science blocker closure ledger and release dossier.\\
Were reviewers addressed? & Reviewer navigation matrix and Tsukrov repair packet.\\
Are domain projections honest? & K0--K12 domain projection and replay state.\\
Is release publication lawful now? & LRGEF state, owner approval packet, and no-send manifest.\\
\end{longtable}

\section*{Five Reviewer Routes}
\begin{longtable}{>{\raggedright\arraybackslash}p{0.25\linewidth}>{\raggedright\arraybackslash}p{0.65\linewidth}}
\textbf{Route} & \textbf{What to verify}\\
\hline
Owner & Read the owner memo, release policy explainer, blocker closure map, and publication locks.\\
Mathematical & Read the definitions, theorem roadmap, proof machinery, and support map.\\
Domain & Read the K0--K12 projection state, domain replay packet, and DRT strict-salvage packet.\\
Reproducibility & Read the manifest, checksums, data manifest, simulation results, and PDF quality report.\\
Publication & Read LRGEF, Parfit, citation metadata, release notes, and postflight checklist.\\
\end{longtable}

\section*{Proven, Replayed, Future}
\begin{longtable}{>{\raggedright\arraybackslash}p{0.28\linewidth}>{\raggedright\arraybackslash}p{0.58\linewidth}}
\textbf{Bucket} & \textbf{Plain meaning}\\
\hline
Proven or evidence-closed & The claim has a proof route, source audit, terminal closure row, or claim-ledger support.\\
Replayed & The claim has a fixed reproducibility route, output hash, tolerance, or domain replay state.\\
Future work & The idea may be important, but 1.3.2 does not ask the reader to accept it as settled.\\
\end{longtable}

\section*{Why the Release Is Stronger Than a Draft}
A draft can hide its weak points. This release cannot. The release machine scans visible PDFs, metadata, package composition, no-send locks, Parfitian governance, source freshness, research packets, and the 82-row terminality ledger. A statement becomes release-facing only when it survives that route.

\section*{What a Careful Reader Should Not Infer}
Do not infer that OC 1.3.2 claims unbounded predictive authority, domain replacement, publication permission, or authority over unrelated problem classes. The package is strong because it states exactly what is proven, replayed, benchmarked, source-audited, and still outside release authority.

\clearpage
\section*{Reader Appendix A: The Whole Release in One Table}
\begin{longtable}{>{\raggedright\arraybackslash}p{0.26\linewidth}>{\raggedright\arraybackslash}p{0.62\linewidth}}
\textbf{Part} & \textbf{Why it matters}\\
\hline
Master monograph & Full theory narrative, dedication, definitions, theorem routes, and appendices.\\
Journal core & Compact scientific kernel for reviewers who need the shortest rigorous route.\\
Readable overview & Plain-language entry point and release navigation.\\
Methods companion & How evidence, replay, manifests, and package gates connect.\\
Critique map & Where objections are answered and how to inspect the replies.\\
Technical spine & Expert navigation through formal support classes and release invariants.\\
\end{longtable}

\clearpage
\section*{Reader Appendix B: How to Avoid Overreading}
If a section says structural support, read it as structural support. If a section says replay, inspect the replay route. If a section says future work, do not treat it as a result of this release. This discipline is what lets the package include a large frontier without pretending that the frontier is already canonical.

\section*{Reader Appendix C: Publication Decision}
The practical owner decision is not whether the release is perfect in the metaphysical sense. The decision is whether this package is honest, current, reproducible, bounded, and worth publishing as the stable 1.3.2 baseline after the no-send locks are deliberately released.

\clearpage
\section*{Reader Appendix D: What Comes After 1.3.2}
After owner review, the scientific queue moves to OC 1.4 formal hardening, math/metamodel strengthening, domain bridge expansion, benchmark growth, and publication-family preparation. Those tasks should cite 1.3.2 as the stable baseline, not silently rewrite it.

\end{document}
